% ===========================================================================
\chapter{Pattern review}
\label{ch:review}

\section{The short answer}

\begin{keyidea}
The structure is right and the wiring is right. What is wrong is concentrated in one place --- the
\code{Repository} trait --- and it has now produced hard evidence against itself. Everything else on
the list below is a small, local fix.
\end{keyidea}

Since the first draft of this document you fixed both blockers it identified in the wiring: the pool
is now built once above \code{HttpServer::new}, and the \code{ssr} gates have moved off the subject
modules onto the files that actually need them. Both builds are green, and there is a working
end-to-end path from Postgres to a button in the browser. That is the hard part of this stack and it
is done.

\section{What is working, and worth not second-guessing}

\begin{table}[htbp]
\centering
\caption{Patterns in the tree that are correct. Listed because knowing which decisions are settled is
as useful as knowing which are not.}
\label{tab:good}
\small
\begin{tabular}{@{}L{0.30\textwidth}L{0.64\textwidth}@{}}
\toprule
\textbf{Pattern} & \textbf{Why it is right} \\
\midrule
Per-file \code{cfg} gates in \file{<subject>/mod.rs} & \code{dto} and \code{components} ungated,
\code{model} and \code{repository} gated. Exactly the shape Ch.~\ref{ch:layers} argues for, and it
is what makes a server function expressible. \\
Pool built once, \code{AppState} cloned per worker & Twelve pools became one.
\code{app\_state.clone()} is an \code{Arc} bump. \\
\code{AppState} holding repositories & Server functions ask for a capability rather than a database
handle. \S\ref{sec:appstate}. \\
Three structs per subject: \code{Item} / \code{Row} / \code{Dto} &
\code{PublicationItem} is \code{Queryable}, \code{PublicationItemRow} is \code{Insertable} with no
\code{id}, \code{PublicationItemDto} is the wire shape. This is \emph{better} than the two-struct
rule in Ch.~\ref{ch:dto} --- a separate insert struct is how you stop passing a meaningless
\code{id: 0} into a create. \\
\code{From<PublicationItemDto> for PublicationItemRow} & Conversion by trait, next to the model, in
the direction the standard library expects. \\
\code{Action} $+$ \code{.into\_any()} in \code{ServerButton} & Correct 0.8 idioms: an \code{Action}
for the imperative call, \code{pending()}/\code{value()} for state, \code{.into\_any()} to reconcile
three branches with different view types. Ch.~\ref{ch:loading}. \\
\code{entrypoint/} as the process edge & The instinct that ``the server lives here'' is right.
Ch.~\ref{ch:routing} fills it in. \\
\bottomrule
\end{tabular}
\end{table}

\section{The one structural problem: \code{Repository<T, U>}}

Chapter~\ref{ch:blocking} argued that this trait would not survive. It has now failed in the exact
way predicted, and the evidence is in the tree rather than in an argument.

\begin{verified}[Verified: the trait is forcing three implementors to lie]
Adding \code{create\_article} to \code{Repository<T, U>} obliged every implementor to provide it.
Three of the four cannot:
\begin{lstlisting}[style=rs]
// src/portfolio/repository.rs
fn create_article(&mut self, article: PortfolioItemDto) -> anyhow::Result<()> { todo!() }

// src/job_experience/repository.rs
fn create_article(&mut self, article: JobExperienceItemDto) -> anyhow::Result<()> { todo!() }

// src/personal_information/repository.rs
fn create_article(&mut self, article: PersonalInformationDto) -> anyhow::Result<()> { todo!() }
\end{lstlisting}
And the compiler is already reporting it, three times over:
\begin{lstlisting}[style=out]
$ cargo check --no-default-features --features ssr
warning: unused variable: `article`
warning: unused variable: `article`
warning: unused variable: `article`
\end{lstlisting}
\end{verified}

Two separate things went wrong, and they are worth naming separately because only one of them is
about the trait.

\textbf{The method name is subject-specific.} ``Article'' is publication vocabulary. A portfolio item
is not an article; neither is a job. The name could not generalise because the \emph{concept} does
not: what a shared trait can express is \code{create}, and even that is only meaningful when every
implementor genuinely creates one thing from one DTO.

\textbf{The trait has no shared behaviour left to describe.} \code{get\_all} and \code{get\_by\_id}
happened to fit all four subjects when they were the only two methods. The third method fits one.
The fourth --- \code{list\_public()} for portfolio, \code{newest\_first()} for publications,
\code{by\_slug()} for both --- will fit fewer. There is no interface here, only a coincidence that has
now expired. Three \code{todo!()}s is what a coincidence looks like when you build on it.

\begin{pattern}[Delete the trait; keep the structs]
This is a smaller change than it sounds, because the bodies already exist and do not move.
\begin{lstlisting}[style=rs,caption={\file{src/publications/repository.rs} --- inherent methods, no trait.}]
impl PublicationsRepository {
    pub fn new(pool: DbPool) -> Self { Self { pool } }

    // &self, not &mut self -- see below
    pub async fn list_newest_first(&self) -> Result<Vec<PublicationItem>, RepoError> { /* ... */ }
    pub async fn by_id(&self, id: i32) -> Result<PublicationItem, RepoError> { /* ... */ }
    pub async fn upsert_many(&self, items: &[PublicationItemDto]) -> Result<usize, RepoError> { /* ... */ }
}
\end{lstlisting}
Delete \file{src/core/mod.rs}, delete the three \code{todo!()} stubs, and delete
\code{use crate::core::Repository;} from the four repositories and from
\file{bin/seed\_publications.rs}. Nothing else changes: \code{AppState} still holds the same four
types, and the call sites keep working because they were calling methods, not the trait.

Reintroduce a trait only when you have a \emph{second implementation} of the same behaviour --- an
in-memory fake for tests is the usual first one. A trait with one implementor is a trait you are
maintaining for free.
\end{pattern}

\begin{why}[But isn't a repository trait ``good practice''?]
It is, in languages where the alternative is worse. The pattern comes from ecosystems where you need
an interface to get dependency injection, mocking, or a runtime-swappable backend. Rust gives you
the first two through ordinary generics and \code{cfg(test)}, and you do not need the third: there
is one Postgres.

Rust's own guidance points the same way --- generics exist to make a function usable by more callers,
not to make a type substitutable in anticipation. And coherence means a generic
\code{Repository<T>} cannot be specialised per entity anyway, so you write the same number of method
bodies with worse type inference and worse error messages. The trait was costing you and paying
nothing, which is why it grew a method that three implementors had to refuse.
\end{why}

\section{Five local fixes, in the order they will bite}

\subsection{1. \code{\&mut self}, and the clone it is forcing}

The cost is now visible in your own code:

\begin{lstlisting}[style=rs,caption={\file{src/publications/server.rs}}]
let state: Data<AppState> = extract().await?;
let mut repo = state.publications_repository.clone();   // <-- cloned only to get `mut`
let results = repo.get_all().unwrap();
\end{lstlisting}

That \code{.clone()} is not there for a reason of its own --- it is there because \code{get\_all}
takes \code{\&mut self} and \code{Data<AppState>} hands out shared references. The trait forces a
clone at every call site to obtain a mutability nobody needs: \code{Pool::get} takes \code{\&self},
and no method in any repository mutates anything. Change the receivers to \code{\&self} and the line
becomes:

\begin{lstlisting}[style=rs]
let results = state.publications_repository.list_newest_first().await?;
\end{lstlisting}

Defect \dnum{8}, and it goes away as part of deleting the trait.

\subsection{2. \code{.unwrap()} in a server function}

\code{repo.get\_all().unwrap()} panics the Actix worker thread when the database is unreachable ---
which is a 500 with no body, no log line, and a poisoned worker, rather than an error the
\code{ErrorBoundary} could render. In a server function, \code{?} is always available:
\code{Result<T, ServerFnError>} has a blanket \code{From} for any \code{std::error::Error}. Use it.

\subsection{3. Diesel is still being called on the async runtime}

\code{server.rs} calls \code{repo.get\_all()} directly inside an \code{async fn}. That is the
blocking bug from Chapter~\ref{ch:blocking}, now live: the query holds an Actix worker thread for its
whole duration. With twelve workers it is invisible; it is invisible right up until it is not. Wrap
the query in \code{web::block} inside the repository, with \code{pool.get()} inside the closure.

\subsection{4. The DTO has no serde derives --- and this is the next wall}

\begin{lstlisting}[style=rs,caption={\file{src/publications/dto.rs}}]
#[derive(Debug, Clone)]          // <-- no Serialize, no Deserialize
pub struct PublicationItemDto { /* ... */ }
\end{lstlisting}

This compiles today only because \code{get\_all\_publications} returns \code{Vec<String>}. The moment
it returns \code{Vec<PublicationItemDto>} --- which is the next thing you will want, since
\code{Vec<String>} throws away the year, journal and link the page needs --- you will hit two errors
in sequence:

\begin{enumerate}
  \item \code{the trait bound `PublicationItemDto: serde::Serialize` is not satisfied}, four times
    per DTO, with notes mentioning \code{BrowserMockRes} and \code{JsonSerdeCodec}. Add
    \code{Serialize, Deserialize} to the derive.
  \item \code{error[E0432]: unresolved import serde} in the \emph{wasm} build --- because
    \code{serde} is still \code{optional = true} with \code{dep:serde} inside the \code{ssr} feature.
    Ch.~\ref{ch:boundary} has the two-line \file{Cargo.toml} fix.
\end{enumerate}

Defect \dnum{6}. Fixing it before you change the return type turns a confusing twenty-minute detour
into a non-event.

\subsection{5. \code{impl Into<...>} instead of \code{impl From<...>}}

\file{bin/seed\_publications.rs} opens with \code{\#![allow(clippy::from\_over\_into)]}, which is the
tell. Clippy is right: implement \code{From<Article> for PublicationItemDto} and you get \code{Into}
generated for free, in both directions of inference, and the \code{allow} disappears. The reason the
lint exists is that a hand-written \code{Into} does not give you the reciprocal \code{From}, so
generic code bounded on \code{From} silently will not accept your type.

\section{Two smaller things while you are in the files}

\begin{itemize}
  \item \textbf{\code{abs} is still \code{abs}.} It is now in three places --- the column, the model,
    and the DTO --- so the rename costs three edits today and more next week. Defect \dnum{15}.
  \item \textbf{\code{components.rs} has a redundant import.}
    \code{use leptos::prelude::*;} followed by \code{use leptos::\{component, view, IntoView\};} ---
    the prelude already provides all three. Harmless, and worth deleting so the next file you write
    copies the right pattern.
\end{itemize}

\section{The revised next step}

Chapter~\ref{ch:roadmap}'s M1 was ``get one slice from Postgres to the screen''. You have done that
in its rough form, so the remaining work in M1 is to make the same path carry real data:

\begin{runbook}[What to do next, concretely]
\begin{enumerate}
  \item Make \code{serde} unconditional in \file{Cargo.toml}, and add
    \code{Serialize, Deserialize} to all four DTOs. \dnum{6}
  \item Delete \file{src/core/mod.rs} and the trait; convert the four repositories to inherent
    methods taking \code{\&self}; delete the three \code{todo!()} stubs. \dnum{8}
  \item Wrap the queries in \code{web::block}, with \code{pool.get()} inside the closure.
  \item Change \code{get\_all\_publications} to return \code{Result<Vec<PublicationItemDto>,
    ServerFnError>}, with \code{?} instead of \code{.unwrap()}.
  \item Replace \code{ServerButton} with a \code{Resource} $+$ \code{Transition} so the list is
    server-rendered rather than fetched on click. Keep the button somewhere until then --- it is a
    genuinely useful hydration test (\S\ref{ch:runcontract}).
  \item Add \code{external\_id} $+$ \code{UNIQUE}, and make the seeder an upsert, before you run it
    again. \S\ref{sec:idempotent}
  \item Extract \file{openalex.rs} and \file{service.rs}; reduce the binary to composition.
    Ch.~\ref{ch:services}
\end{enumerate}
\end{runbook}

Steps 1--4 are perhaps an hour and remove four of the five defects in this chapter. Step 6 is the one
with a deadline attached: it needs to happen before the next \code{cargo run --bin
seed\_publications}, or you will be deduplicating rows by hand.
